\chapter{Options and Derivatives}\label{ch:iv:options-and-derivatives}

A screen shows a call at 7.10 and a put at 2.40 on the same stock, with the same strike of 100 and the same
three-month expiry; the stock is at 104, it pays no dividend before expiry, and rates are 4\%. The candidate has ten
seconds to say whether there is money on the table and, if so, what to trade. Option questions are the core of
market-maker and bank interviews, and they test three layers: the no-arbitrage relations that hold whatever the
model, the Greeks as intuition rather than formulas, and what a hedged book earns and loses. The theory is in One
Quant Book~1, chapters 23 to~26, and One Quant Book~5, chapters 1 to~15.

\section{No-arbitrage: parity, bounds and convexity}

Three model-free facts cover most questions. \emph{Parity}: for European options on a stock with dividends of
present value $D$, $C - P = S - D - Ke^{-rT}$. \emph{Bounds}: $\max(S - D - Ke^{-rT}, 0) \le C \le S$, and $\max(Ke^{-rT} - S
+ D, 0) \le P \le Ke^{-rT}$. \emph{Shape in the strike}: call prices fall with the strike, by no more than the discounted
strike difference, and are convex, so every butterfly has a non-negative price (One Quant Book~5, chapter~1).

\begin{method}[Checking quotes for an arbitrage]\label{met:iv:options-and-derivatives:arb}
\begin{enumerate}
\item Parity first: compute $S - D - Ke^{-rT}$ and compare with the quoted $C - P$.
\item Then the bounds for each option on its own.
\item Then the strike structure: monotone, slopes bounded, convex (butterflies and call spreads).
\item For each violation, write the trade (buy the cheap side, sell the rich side, hedge with stock and cash) and its
locked-in profit; mention what could make the ``arbitrage'' illusory: an unmodelled dividend, a hard-to-borrow stock,
American exercise, stale quotes.
\end{enumerate}
\end{method}

\begin{example}[The hook's quotes]\label{ex:iv:options-and-derivatives:hook}
$S - Ke^{-rT} = 104 - 100e^{-0.01} \approx 4.995$, and the quoted $C - P = 4.70$: the call is cheap relative to the put by
about 0.295. Buy the call, sell the put, short the stock and lend $100e^{-0.01}$: the position is worth zero at expiry
whatever happens, and it costs $-0.295$ today. Before trading, ask about dividends and borrow: a dividend of present value
0.30 or a borrow fee of about 1.1\% a year would explain the gap.
\end{example}

\section{Greeks by intuition}

Interviewers ask for signs, shapes and peaks more often than for formulas. A long call has positive delta, gamma and
vega, and negative theta; a long put has negative delta and the same signs of gamma and vega. Gamma and theta of an
at-the-money option grow as expiry approaches, like $1/\sqrt T$; vega shrinks like $\sqrt T$
(\cref{fig:iv:options-and-derivatives:gamma}). The at-the-money approximation $C \approx 0.4\,S\sigma\sqrt T$ (from
$\varphi(0) \approx 0.4$) prices an at-the-money option in the head.

\begin{omfigure}
\begin{tikzpicture}
\begin{axis}[width=10cm,height=5.4cm,xlabel={\small spot},ylabel={\small gamma},
  xmin=80,xmax=120,ymin=0,ymax=0.15,
  yticklabel style={/pgf/number format/fixed,/pgf/number format/precision=2},scaled y ticks=false,
  tick label style={font=\footnotesize},grid=major,grid style={gray!20},table/col sep=comma,
  legend cell align=left,legend style={font=\scriptsize,at={(0.5,-0.32)},anchor=north,
  /tikz/every even column/.append style={column sep=8pt}},legend columns=3]
\addplot[omThm,thick] table[x=spot,y=week]{figdata/interviews/17-options-and-derivatives/gamma.csv};
\addplot[omProp,thick,dashed] table[x=spot,y=quarter]{figdata/interviews/17-options-and-derivatives/gamma.csv};
\addplot[omDef,thick,densely dotted] table[x=spot,y=year]{figdata/interviews/17-options-and-derivatives/gamma.csv};
\legend{one week,three months,one year}
\end{axis}
\end{tikzpicture}

\omcaption{Black--Scholes gamma of an option struck at 100 against the spot, at 20\% volatility and zero rates, for
three times to expiry. Near expiry gamma is concentrated in a narrow band around the strike and is about seven times
its one-year value at the money. Data: \texttt{fig\_iv\_greeks.py}.}\label{fig:iv:options-and-derivatives:gamma}
\end{omfigure}

\begin{example}[Gamma and theta of a hedged book]\label{ex:iv:options-and-derivatives:pnl}
A delta-hedged book earns $\tfrac12\Gamma(\Delta S)^2$ from a move and pays theta, which at the implied volatility is
$\tfrac12\Gamma S^2\sigma^2\,\Delta t$ (One Quant Book~5, chapter~4). With a position gamma of 50 (shares per unit of
spot), the stock at 100 and implied volatility 20\%, a day's theta is about $-39.7$; a move of 2 earns 100, a net of about
60. The break-even move is $S\sigma/\sqrt{252} \approx 1.26$: the hedged book is long realised against implied
volatility.
\end{example}

\section{The smile and what it implies}

Call prices across strikes encode the market's risk-neutral distribution. The digital call is minus the strike
derivative of the call, and the density is the second derivative (Breeden and Litzenberger, 1978). On a smile, the strike
derivative of the call includes the change of implied volatility with the strike, so the digital's price is $\Phi(d_2)$
minus vega times the skew's slope $\partial\sigma/\partial K$.

\begin{example}[A digital under skew]\label{ex:iv:options-and-derivatives:digital}
At $S = K = 100$, $\sigma = 20\%$, one year and zero rates, the flat-volatility digital is worth $\Phi(-0.1) \approx 0.460$
and vega is about 39.7. With a skew of $-0.1$ volatility point per unit of strike ($\partial\sigma/\partial K = -0.001$), the
digital is worth $0.460 + 39.7 \times 0.001 \approx 0.500$: a negative skew makes upside digitals more expensive than flat
volatility says, because the calls just above the strike are relatively cheap.
\end{example}

\section{Hedging scenarios}

Scenario questions combine the pieces: a short straddle before an event, a book of zero-day options into the close, a
position that is flat in delta but not in gamma. The answers share a structure: say what the position is long or short
(delta, gamma, vega, the skew), compute the P\&L of the move with $\tfrac12\Gamma(\Delta S)^2$ and the vega change, say
what hedging would cost, and name the risk that the Greeks miss (a jump, pin risk at expiry, early exercise, a borrow
recall).

\section{Worked answers}

\begin{example}[One period, two states]\label{ex:iv:options-and-derivatives:binomial}
\emph{``A stock at 100 will be at 120 or at 90 tomorrow; rates are zero. Price the call struck at 100.''} Replicate: the
call pays 20 or 0, so hold $\Delta = (20 - 0)/(120 - 90) = 2/3$ of a share and borrow what makes the portfolio pay 0 in the
down state, $\tfrac23 \times 90 = 60$. The portfolio costs $\tfrac23 \times 100 - 60 = 6.67$, and so does the call. The
same number comes from the probability that makes the stock a martingale, $q = (100 - 90)/(120 - 90) = 1/3$: $20q = 6.67$.
The real probability of the up move never entered; the interviewer asks next ``what if it is 90\%?'', and the answer is
that the price does not change, because the hedge does not depend on it.
\end{example}

\begin{example}[A risk reversal when the skew steepens]\label{ex:iv:options-and-derivatives:rr}
\emph{``You are long a risk reversal: long an out-of-the-money call, short an out-of-the-money put, each with a vega of 0.30
per volatility point. The put's implied volatility rises 2 points and the call's falls 1. What is your P\&L, ignoring the
spot move?''} The call loses $1 \times 0.30 = 0.30$; the short put loses $2 \times 0.30 = 0.60$; total $-0.90$. The position
was nearly flat vega, and it still lost, because it was short skew: a vega number that nets calls against puts hides the
exposure to the difference between their volatilities, and an interviewer who asks this is checking that the candidate
knows it.
\end{example}

\begin{example}[Theta at the money, in one line]\label{ex:iv:options-and-derivatives:theta}
\emph{``What does a three-month at-the-money call on a stock at 100 with 20\% volatility lose a day, with zero rates?''} At
the money with zero rates, the price is about $0.4\,\sigma S\sqrt{T}$, and theta is its time derivative, $-\sigma
S/(2\sqrt{2\pi T})$: here $-20/(2 \times 1.2533) \approx -7.98$ a year, about 0.022 a calendar day. \emph{Check} against
gamma: for a hedged book, theta pays for gamma, $\Theta = -\tfrac12\Gamma S^2\sigma^2$, and the at-the-money gamma
$\varphi(0.05)/(S\sigma\sqrt T) \approx 0.0398$ gives $-\tfrac12 \times 0.0398 \times 10^4 \times 0.04 \approx -7.97$. The
two routes agree to the accuracy of the first, which used $\varphi(0)$ for $\varphi(0.05)$.
\end{example}

\section{Question bank}

\begin{interviewq}[$\star$ \iqroles{trader} \iqfirm{market maker}]\label{iq:iv:options-and-derivatives:1}
The hook's quotes: call 7.10, put 2.40, strike 100, three months, stock 104, rates 4\%, no dividend. Is there an arbitrage?
What do you trade and what do you lock in? What would you check first?
\end{interviewq}

\begin{interviewq}[$\star$ \iqroles{trader} \iqfirm{market maker}]\label{iq:iv:options-and-derivatives:2}
A three-month European call struck at 40 on a non-dividend stock at 50 is quoted at 9.50; rates are 4\%. What is wrong?
\end{interviewq}

\begin{interviewq}[$\star$ \iqroles{trader, bank} \iqfirm{bank}]\label{iq:iv:options-and-derivatives:3}
Give the signs of delta, gamma, vega and theta for a long put. Is there a case in which the theta of a long European put
is positive?
\end{interviewq}

\begin{interviewq}[$\star$ \iqroles{trader} \iqfirm{market maker}]\label{iq:iv:options-and-derivatives:4}
Price a three-month at-the-money call on a stock at 100 with 20\% volatility and zero rates, in your head.
\end{interviewq}

\begin{interviewq}[$\star$ \iqroles{trader} \iqfirm{market maker}]\label{iq:iv:options-and-derivatives:5}
Calls struck at 90, 100 and 110 are quoted at 12, 7 and 1.50. Is anything wrong? What do you trade?
\end{interviewq}

\begin{interviewq}[$\star\star$ \iqroles{trader, bank} \iqfirm{bank}]\label{iq:iv:options-and-derivatives:6}
Why is an American call on a non-dividend stock never exercised early? Give an American put that should be exercised at
once, with numbers.
\end{interviewq}

\begin{interviewq}[$\star\star$ \iqroles{trader, risk} \iqfirm{market maker}]\label{iq:iv:options-and-derivatives:7}
Compare an at-the-money option's gamma at one week and at one year to expiry (spot 100, volatility 20\%). By
what factor do they differ, and why?
\end{interviewq}

\begin{interviewq}[$\star\star$ \iqroles{trader} \iqfirm{market maker}]\label{iq:iv:options-and-derivatives:8}
Your delta-hedged book has a gamma of 50 (shares per unit of spot); the stock is at 100 and implied volatility is 20\%. The
stock moves by 2 today. What is your P\&L, net of theta? What move makes you break even?
\end{interviewq}

\begin{interviewq}[$\star\star$ \iqroles{trader, bank} \iqfirm{bank}]\label{iq:iv:options-and-derivatives:9}
What is the sign of the vega of a cash-or-nothing digital call that is out of the money? In the money? Explain without
formulas.
\end{interviewq}

\begin{interviewq}[$\star\star$ \iqroles{bank, trader} \iqfirm{bank}]\label{iq:iv:options-and-derivatives:10}
At-the-money volatility is 20\% and the skew is $-0.1$ volatility point per unit of strike. Price a one-year at-the-money
digital call (spot 100, zero rates), and say why the skew matters.
\end{interviewq}

\begin{interviewq}[$\star\star\star$ \iqroles{researcher, bank} \iqfirm{bank}]\label{iq:iv:options-and-derivatives:11}
Calls with the same expiry struck at 95, 100 and 105 are worth 9.0, 6.0 and 3.8 (zero rates). What is the risk-neutral
probability that the stock ends between 97.5 and 102.5, approximately?
\end{interviewq}

\begin{interviewq}[$\star\star\star$ \iqroles{trader, bank} \iqfirm{market maker}]\label{iq:iv:options-and-derivatives:12}
Why does a variance swap's fair strike exceed at-the-money implied volatility when the skew is negative? With a smile
$\sigma(K) = 20\% - 10\% \ln(K/100)$ (floored at 5\%), one year and zero rates, estimate the fair volatility strike.
\end{interviewq}

\begin{interviewq}[$\star\star\star$ \iqroles{trader} \iqfirm{market maker}]\label{iq:iv:options-and-derivatives:13}
A stock at 100 pays a dividend of 2 in three months; rates are zero. What does put--call parity give for six-month
European options struck at 100? What changes if the options are American?
\end{interviewq}

\begin{interviewq}[$\star\star\star$ \iqroles{trader, risk} \iqfirm{proprietary firm}]\label{iq:iv:options-and-derivatives:14}
You are short an at-the-money straddle with one day to expiry on a stock at 100 with 20\% implied volatility. What did
you receive, what is your break-even move, and what do you lose if the stock moves 3\% by the close? What risk do the
Greeks miss here?
\end{interviewq}

\begin{omsources}
\item D.~T.~Breeden and R.~H.~Litzenberger, ``Prices of state-contingent claims implicit in option prices'',
\emph{Journal of Business} 51(4), 1978.
\item One Quant Book~1, chapters 23--26; One Quant Book~5, chapters 1--6, 14--15 (no-arbitrage, Black--Scholes, Greeks,
American options, variance swaps, digitals).
\end{omsources}
